\clearpage
\section{The clock, and the display it feeds}

\subsection{Two clocks behind one type name}

\file{timer.rs} defines a private \texttt{Clock} enum with one variant per platform.
This is the only piece of the core that is platform-conditional, and it is 40 lines.

\begin{lstlisting}[caption={\file{timer.rs} lines 4--59, abridged.}]
#[derive(Clone, Copy, Debug)]
struct Clock {
    #[cfg(target_os = "macos")]      ticks:  u64,              // mach_continuous_time
    #[cfg(not(target_os = "macos"))] instant: std::time::Instant,
}
impl Clock {
    fn now() -> Self { /* one branch per platform */ }
    fn elapsed(self) -> Duration {
        #[cfg(target_os = "macos")]
        {
            static SCALE: OnceLock<(u32,u32)> = OnceLock::new();
            let &(numer, denom) = SCALE.get_or_init(|| { /* mach_timebase_info */ });
            let nanos = (Clock::now().ticks - self.ticks) as u128 * numer as u128
                        / denom as u128;
            Duration::from_nanos(nanos.min(u64::MAX as u128) as u64)
        }
        #[cfg(not(target_os = "macos"))]
        { self.instant.elapsed() }
    }
}
unsafe extern "C" { fn mach_continuous_time() -> u64; fn mach_timebase_info(..) -> i32; }
\end{lstlisting}

\begin{figure}[H]
\centering
\begin{tikzpicture}[
  c/.style={rounded corners=2pt,draw=ink!50,fill=white,align=center,inner sep=5pt,
            font=\scriptsize\sffamily,text width=38mm,minimum height=12mm},
  mac/.style={c,draw=mint!70!black,fill=mint!18},
  oth/.style={c,draw=pastel!45!black,fill=corefill},
]
  \node[c, draw=softink!50, fill=soft, text width=44mm] (api) at (0,0)
    {\textbf{Clock::now()} and \texttt{Clock::elapsed()}\\[1mm]
     {\tiny the rest of \texttt{timer.rs} never asks which platform it is on}};

  \node[mac] (m) at (-5.2,-3.4) {\textbf{macOS}\\[1mm]\texttt{mach\_continuous\_time()}\\[1mm]
     {\tiny raw ticks, scaled by\\\texttt{mach\_timebase\_info()}\\\textit{numer/denom}}};
  \node[oth] (o) at (5.2,-3.4) {\textbf{Linux / Windows}\\[1mm]\texttt{std::time::Instant::now()}\\[1mm]
     {\tiny the monotonic clock\\\textit{that the OS guarantees}}};
  \draw[e] (api.south) -- (m.north);
  \draw[e] (api.south) -- (o.north);
\end{tikzpicture}
\caption{\texttt{Clock} is the entire platform abstraction of the core.}
\label{fig:clock}
\end{figure}

\keybox{mint}{\textbf{Why \code{mach\_continuous\_time} and not
\code{mach\_absolute\_time}?} The ``continuous'' variant keeps counting while the
machine is asleep, which is what a Pomodoro timer wants: if the lid is shut for ten
minutes, the session really did advance by ten minutes and
\texttt{App::sleep()} decides separately whether to pause it. The timebase
conversion is cached in a \code{OnceLock} because it is a constant for the lifetime
of the process.}

\subsection{Rounding: why a fresh timer shows 25:00, not 24:59}

Two tiny functions carry all of the display logic, and both exist because the
alternative was a visible one-second jitter.

\begin{lstlisting}[caption={Round up, so the label is stable for a whole second.},basicstyle=\ttfamily\footnotesize]
pub fn remaining\_display(\&self) -> Duration {
    let remaining = self.remaining();
    if remaining.is_zero() { return Duration::ZERO; }
    Duration::from_secs(remaining.as_secs() + u64::from(remaining.subsec_nanos() != 0))
}
\end{lstlisting}

\begin{lstlisting}[caption={Sleep exactly until the label would change --- not,basicstyle=\ttfamily\footnotesize],basicstyle=\ttfamily\footnotesize],basicstyle=\ttfamily\footnotesize],basicstyle=\ttfamily\footnotesize],basicstyle=\ttfamily\footnotesize],basicstyle=\ttfamily\footnotesize],basicstyle=\ttfamily\footnotesize],basicstyle=\ttfamily\footnotesize],basicstyle=\ttfamily\footnotesize],basicstyle=\ttfamily\footnotesize],basicstyle=\ttfamily\footnotesize],basicstyle=\ttfamily\footnotesize],basicstyle=\ttfamily\footnotesize],basicstyle=\ttfamily\footnotesize],basicstyle=\ttfamily\footnotesize],basicstyle=\ttfamily\footnotesize],basicstyle=\ttfamily\footnotesize],basicstyle=\ttfamily\footnotesize],basicstyle=\ttfamily\footnotesize],basicstyle=\ttfamily\footnotesize],basicstyle=\ttfamily\footnotesize],basicstyle=\ttfamily\footnotesize],basicstyle=\ttfamily\footnotesize],basicstyle=\ttfamily\footnotesize],basicstyle=\ttfamily\footnotesize],basicstyle=\ttfamily\footnotesize],basicstyle=\ttfamily\footnotesize],basicstyle=\ttfamily\footnotesize],basicstyle=\ttfamily\footnotesize],basicstyle=\ttfamily\footnotesize],basicstyle=\ttfamily\footnotesize],basicstyle=\ttfamily\footnotesize],basicstyle=\ttfamily\footnotesize],basicstyle=\ttfamily\footnotesize],basicstyle=\ttfamily\footnotesize],basicstyle=\ttfamily\footnotesize],basicstyle=\ttfamily\footnotesize],basicstyle=\ttfamily\footnotesize],basicstyle=\ttfamily\footnotesize],basicstyle=\ttfamily\footnotesize],basicstyle=\ttfamily\footnotesize],basicstyle=\ttfamily\footnotesize],basicstyle=\ttfamily\footnotesize],basicstyle=\ttfamily\footnotesize],basicstyle=\ttfamily\footnotesize],basicstyle=\ttfamily\footnotesize],basicstyle=\ttfamily\footnotesize],basicstyle=\ttfamily\footnotesize],basicstyle=\ttfamily\footnotesize],basicstyle=\ttfamily\footnotesize],basicstyle=\ttfamily\footnotesize],basicstyle=\ttfamily\footnotesize],basicstyle=\ttfamily\footnotesize],basicstyle=\ttfamily\footnotesize],basicstyle=\ttfamily\footnotesize],basicstyle=\ttfamily\footnotesize],basicstyle=\ttfamily\footnotesize],basicstyle=\ttfamily\footnotesize],basicstyle=\ttfamily\footnotesize],basicstyle=\ttfamily\footnotesize],basicstyle=\ttfamily\footnotesize],basicstyle=\ttfamily\footnotesize],basicstyle=\ttfamily\footnotesize],basicstyle=\ttfamily\footnotesize],basicstyle=\ttfamily\footnotesize],basicstyle=\ttfamily\footnotesize],basicstyle=\ttfamily\footnotesize],basicstyle=\ttfamily\footnotesize],basicstyle=\ttfamily\footnotesize],basicstyle=\ttfamily\footnotesize],basicstyle=\ttfamily\footnotesize],basicstyle=\ttfamily\footnotesize],basicstyle=\ttfamily\footnotesize],basicstyle=\ttfamily\footnotesize],basicstyle=\ttfamily\footnotesize],basicstyle=\ttfamily\footnotesize],basicstyle=\ttfamily\footnotesize],basicstyle=\ttfamily\footnotesize],basicstyle=\ttfamily\footnotesize],basicstyle=\ttfamily\footnotesize],basicstyle=\ttfamily\footnotesize],basicstyle=\ttfamily\footnotesize],basicstyle=\ttfamily\footnotesize],basicstyle=\ttfamily\footnotesize],basicstyle=\ttfamily\footnotesize],basicstyle=\ttfamily\footnotesize],basicstyle=\ttfamily\footnotesize],basicstyle=\ttfamily\footnotesize],basicstyle=\ttfamily\footnotesize],basicstyle=\ttfamily\footnotesize],basicstyle=\ttfamily\footnotesize],basicstyle=\ttfamily\footnotesize],basicstyle=\ttfamily\footnotesize],basicstyle=\ttfamily\footnotesize],basicstyle=\ttfamily\footnotesize],basicstyle=\ttfamily\footnotesize],basicstyle=\ttfamily\footnotesize],basicstyle=\ttfamily\footnotesize],basicstyle=\ttfamily\footnotesize],basicstyle=\ttfamily\footnotesize],basicstyle=\ttfamily\footnotesize],basicstyle=\ttfamily\footnotesize],basicstyle=\ttfamily\footnotesize],basicstyle=\ttfamily\footnotesize],basicstyle=\ttfamily\footnotesize],basicstyle=\ttfamily\footnotesize],basicstyle=\ttfamily\footnotesize],basicstyle=\ttfamily\footnotesize],basicstyle=\ttfamily\footnotesize],basicstyle=\ttfamily\footnotesize],basicstyle=\ttfamily\footnotesize],basicstyle=\ttfamily\footnotesize],basicstyle=\ttfamily\footnotesize],basicstyle=\ttfamily\footnotesize],basicstyle=\ttfamily\footnotesize],basicstyle=\ttfamily\footnotesize],basicstyle=\ttfamily\footnotesize],basicstyle=\ttfamily\footnotesize],basicstyle=\ttfamily\footnotesize],basicstyle=\ttfamily\footnotesize],basicstyle=\ttfamily\footnotesize],basicstyle=\ttfamily\footnotesize],basicstyle=\ttfamily\footnotesize],basicstyle=\ttfamily\footnotesize],basicstyle=\ttfamily\footnotesize],basicstyle=\ttfamily\footnotesize],basicstyle=\ttfamily\footnotesize]
\texttt{duration - remaining}, the complement, and busy.}]
pub fn next_display_change(&self) -> Duration {
    let remaining = self.remaining();
    if remaining.is_zero() { return Duration::from_secs(1); }
    let nanos = remaining.subsec_nanos();
    if nanos == 0 { Duration::from_secs(1) } else { Duration::from_nanos(u64::from(nanos)) }
}
\end{lstlisting}

\begin{figure}[H]
\centering
\begin{tikzpicture}[x=1mm,y=1mm]
  \draw[draw=softink!30,line width=0.5pt] (0,0) rectangle (150,62);
  % timeline
  \draw[-stealth, draw=ink!60] (8,46) -- (144,46);
  \node[anchor=south west,font=\tiny\sffamily,text=softink] at (8,47) {\texttt{elapsed}};

  % sub-second slivers
  \foreach \fx/\fw/\flab in {20/16/{24:59.25}, 44/12/{24:59.75}, 68/16/{24:00.25}, 92/12/{24:00.75}, 116/14/{zero}}
  {
    \fill[accent!30] (\fx,42) rectangle ({\fx+\fw},44.5);
    \draw[softink!40] (\fx,42) -- (\fx,47);
    \node[anchor=north,font=\tiny\sffamily,text=softink,align=center] at ({\fx+\fw/2},40) {\flab};
  }
  % labels
  \node[anchor=south,font=\tiny\sffamily\bfseries,text=softink] at (28,53) {displayed value};
  \node[anchor=south,font=\tiny\sffamily\bfseries,text=accent!55!black] at (100,53) {wake-up instants};
  \node[anchor=west,font=\tiny\sffamily,text=accent!55!black] at (116,56) {one per second};
  \draw[er] (112,55) -- (118,47);

  \end{tikzpicture}
\end{figure}

\begin{center}\small
\begin{minipage}{0.92\linewidth}
\texttt{remaining\_display()} rounds \emph{up}: with $0.25$\,s left the label still
prints \texttt{24:59}, so it changes exactly when the fractional part reaches zero.

\texttt{next\_display\_change()} returns the \emph{complement} of that, i.e.\ the time
\emph{until} the boundary --- $0.75$\,s at the first sample above, not $0.25$\,s.
\end{minipage}
\end{center}

\begin{figure}[H]
\centering
\caption{The two functions exist to keep each other honest: round up for the label, and
sleep until the label would change.}
\label{fig:rounding}
\end{figure}

\subsection{How each shell spends that delay}

\begin{table}[H]
\centering\small
\begin{tabular}{@{}p{30mm} p{66mm} p{62mm}@{}}
\toprule
& \textbf{mac\_shell.rs} & \textbf{egui\_shell.rs} \\
\midrule
Idle timer & no timer armed & no repaint scheduled \\
Running & \texttt{refresh()} ends with a one-shot \texttt{NSTimer} &
\texttt{schedule\_repaint()} calls \texttt{ctx.request\_repaint\_after(d)} \\
Tolerance & \texttt{setTolerance(0.02)} --- the OS may fire 20\,ms late &
exact; egui coalesces into the next frame \\
Clamped to & $\max(0.01\,\text{s})$ & n/a \\
Repeat? & \texttt{repeats: false}, then re-armed in \texttt{refresh()} & re-armed every frame \\
\bottomrule
\end{tabular}
\caption{Both shells converge on the same rule: schedule exactly one wake-up, at the
next visible change. Neither one ever busy-waits.}
\end{table}

\begin{lstlisting}[caption={\file{app.rs} lines 172--176: the whole scheduling policy is
three lines.}]
#[cfg(feature = "egui-ui")]
pub fn schedule_repaint(&self, ctx: &egui::Context) {
    if self.timer.state() == RunState::Running {
        ctx.request_repaint_after(self.timer.next_display_change());
    }
}
\end{lstlisting}

\begin{figure}[H]
\centering
\begin{tikzpicture}[
  b/.style={rounded corners=2pt,draw=ink!50,fill=white,align=center,inner sep=4pt,
            font=\tiny\sffamily,text width=24mm,minimum height=9mm},
  k/.style={b,draw=accent!65,fill=accent!12,text=accent!55!black},
  g/.style={b,draw=mint!70!black,fill=mint!18,text=mint!60!black},
]
  \node[b] (a) at (0,0) {\textbf{running}\\{\tiny state == \texttt{Running}}};
  \node[g] (b) at (0,-1.7) {\tiny arm a wake-up at\\\texttt{next\_display\_change()}};
  \node[k] (c) at (0,-3.4) {\tiny wake up, update\\\texttt{last\_second}, re-arm};
  \node[b] (d) at (0,-5.1) {\textbf{idle / paused}\\{\tiny nothing is armed}};
  \draw[es] (a)--(b); \draw[es] (b)--(c);
  \draw[ethin, softink] (c.east) -- ++(9mm,0) |- (a.east) node[pos=0.25,right,font=\tiny,text=softink] {\scriptsize repeats};
  \draw[ethin, softink!40, dashed] (d.east) -- ++(12mm,0) node[right,font=\tiny,text=softink] {\scriptsize 0 events/s};
\end{tikzpicture}
\caption{The repaint loop, which is really just a one-node cycle.}
\label{fig:repaint}
\end{figure}